%% The first command in your LaTeX source must be the \documentclass command.
%%
%% Options:
%% twocolumn : Two column layout.
%% hf: enable header and footer.
\documentclass[
% twocolumn,
% hf,
]{ceurart}

%%
%% One can fix some overfulls
\sloppy
\usepackage{float}
%%
%% Minted listings support 
%% Need pygment <http://pygments.org/> <http://pypi.python.org/pypi/Pygments>
\usepackage{listings}
%% auto break lines
\lstset{breaklines=true}

%%
%% end of the preamble, start of the body of the document source.
\begin{document}

%%
%% Rights management information.
%% CC-BY is default license.
\copyrightyear{2024}
\copyrightclause{Copyright for this paper by its authors.
  Use permitted under Creative Commons License Attribution 4.0
  International (CC BY 4.0).}

%%
%% This command is for the conference information
\conference{ICPM 2024 Tool Demonstration Track, October 14-18, 2024, Kongens Lyngby, Denmark}

%%
%% The "title" command
\title{PowerMine: A Tool for Process Mining}

\tnotemark[1]
\tnotetext[1]{You can use this document as the template for preparing your
  publication.}

%%
%% The "author" command and its associated commands are used to define
%% the authors and their affiliations.
\author[1]{Alice A. Andersen}[%
orcid=0000-0000-0000-0000,
email=andersen@uanywhere.eu,
]
\cormark[1]
\fnmark[1]
\address[1]{University of Anywhere,
  University St. 42, 12345 Happyville, Country}

\author[2]{Bob B. Baron}[%
orcid=9999-9999-9999-9999,
email=b.baron@acme.com,
]
\fnmark[1]
\address[2]{ACME Inc., Profit St. 123, 12345 Happyville, Country}

\fnmark[1]

%% Footnotes
\cortext[1]{Corresponding author.}
\fntext[1]{These authors contributed equally.}

%%
%% The abstract is a short summary of the work to be presented in the
%% article.
\begin{abstract}
  Here goes the abstract.
\end{abstract}

%%
%% Keywords. The author(s) should pick words that accurately describe
%% the work being presented. Separate the keywords with commas.
\begin{keywords}
  Keyword \sep
  Another keyword 
\end{keywords}

%%
%% This command processes the author and affiliation and title
%% information and builds the first part of the formatted document.
\maketitle

%% Tool metadata table: if a metatadum is not applicable for your tool 
%% (e.g., source code repository for a commercial application), write 
%% N/A in the appropriate field
%% Please do not move/resize this table - as it should stay after the 
%% abstract, before the first section - nor change the labels in the 
%% metadata description column
\begin{table*}[h]
\footnotesize
\begin{tabular}{ll}
\toprule
Metadata description               & Value \\
\midrule
Tool name                          & PowerMine \\
Current version                    & 1.0 \\
Legal code license                 & Apache 2.0 \\
Languages, tools and services used & C++, Qt, GraphViz \\
Supported operating environment    & Microsoft Windows, GNU/Linux  \\
Download/Demo URL                  & \url{hxxp://github.com/uanywhere/powermine/releases/tag/v1.0}  \\
Documentation URL                  & \url{hxxp://uanywhere.notion.site/ICPM-536b11fc596f4ba6aa631f0a4b69ed5e} \\
Source code repository             & \url{hxxp://github.com/uanywhere/powermine} \\
Screencast video                   & \url{hxxp://www.youtube.com/watch?v=vvOA\_zcOOh4} \\
\bottomrule
\end{tabular}
\end{table*}

\section{Introduction}

Here goes the introduction

\section{Another Section}

Here goes the other content \cite{katz2009research}, which can be freely organized into sections and subsections.

\section{Conclusion}

Here goes the conclusion

%%
%% The acknowledgments section is defined using the "acknowledgments" environment
%% (and NOT an unnumbered section). This ensures the proper
%% identification of the section in the article metadata, and the
%% consistent spelling of the heading.
\begin{acknowledgments}
This template makes use of the CEURART class, which belongs to CEUR-WS.

If needed, acknowledgements should go here. 
\end{acknowledgments}

%%
%% Define the bibliography file to be used
\bibliography{sample-ceur}

%%
%% If your work has an appendix, this is the place to put it.
\appendix

\section{Getting a license}
If the tool requires a license, this appendix should describe how to obtain a (temporary) license. The procedure to obtain the license must not disclose the identity of the reviewers. This appendix will not be included in the final version for the proceedings, if the demo is accepted.

\end{document}

%%
%% End of file
